\documentclass[titlepage]{ctexart}

\usepackage{amsmath}
\usepackage{amssymb}

\title{第七章~概率}
\author{dovego}
\date{\today}

\begin{document}
\maketitle


\section{随机现象与随机事件}

\subsection{随机现象}

确定性现象:在一定条件下必然发生的现象.

随机现象:在一定条件下进行试验或观察会出现不同结果,并且事先无法预测哪一种结果会出现的现象.

随机现象具有两个特点：
(1)至少有两个结果;
(2)事先不知道会出现哪一个结果.

\subsection{样本空间}

试验:对随机现象进行观察或实验的行为.一般用 $E$ 来表示.所得结果称为\textbf{试验结果}

样本空间:试验 $E$ 的所有可能结果组成的集合,记作 $\varOmega$ .

样本点:样本空间 $\varOmega$ 的元素,即试验 $E$ 的每种可能结果.记作 $\omega$.

有限样本空间:样本点的个数为有限个的样本空间.

\subsection{随机事件}

随机事件:试验 $E$ 的样本空间 $\varOmega$ 的子集叫做 $E$ 的随机事件,通常用 $A,B,C,\dots$ 来表示.
当试验结果出现了随机事件中的任意一个样本点 $\omega$ 时,就说发生了该随机事件;反之,如果说随机事件发生了,也就意味着出现了其中的任意一个样本点.

必然事件: $\varOmega$ 本身是样本空间 $\varOmega$ 的子集,所以它也是一个随机事件.
又因为它包含所有样本点,在每一次试验时必然出现其中一个样本点 $\omega$ ,故称 $\varOmega$ 为必然事件.

不可能事件: $\varnothing$ 是样本空间 $\varOmega$ 的子集,所以它也是一个随机事件.
因为它不包含任何样本点,所以在每一次试验时它都不可能发生,故称 $\varnothing$ 为不可能事件.

\subsection{随机事件的运算}

随机事件的运算其实就是集合间的运算:交(积事件)、并(和事件)、补(对立事件)、不相交(互斥事件).

用Venn图可以形象化理解.

交/积事件:事件 $A$ 与 $B$ 同时发生的事件.记作 $A \cap B$ 或 $AB$.

并/和事件:事件 $A$ 与 $B$ 中至少有一个发生的事件.记作 $A \cup B $ 或 $A + B$.

互斥事件:不可能同时发生的事件.若事件 $A$ 与 $B$ 满足 $A \cap B = \varnothing$,则称 $A$ 与 $B$ 为互斥事件.

对立事件:要么发生一个,要么发生另一个.若事件 $A$ 与 $B$ 满足 $A \cup B = \varOmega$ 且 $A \cap B = \varnothing$,则称 $A$ 与 $B$ 互为对立事件.
事件 $A$ 的对立事件记作 $\overline{A}$.


\section{古典概型}

\subsection{古典概型的概率计算公式}

概率:通常用一个数 $P(A)(0 \leqslant P(A) \leqslant 1)$ 来表示事件发生的可能性大小.
特别地, $P(\varOmega) = 1, P(\varnothing) = 0$.

若试验 $E$ 满足:

(1)有限性:样本空间 $\varOmega$ 是有限样本空间;

(2)等可能性:每次试验中, 各个样本点出现的可能性相等.

则称这样的试验模型为 \textbf{古典概率模型}, 简称 \textbf{古典概型}.

在古典概型中,若样本空间 $\varOmega$ 包含的样本点个数为 $n$, 随机事件 $A$ 包含的样本点个数为 $m$,则事件 $A$ 发生的概率为:
\[
P(A) = \frac
    {A\text{包含的样本点个数}}
    {\varOmega \text{包含的样本点个数}}
     = \frac{m}{n}
\]

\subsection{古典概型的应用}

例题 3 的解法 4 相对不太好理解. 这个例子说明抽签、摸球的中奖可能性大小与顺序无关!!! (以前的教材有这个例题吗?解决了我多年的疑问!!!)

互斥事件的概率加法公式:
\[
P(A \cup B) = P(A) + P(B)
\]

\begin{itemize}
    \item (有限可加性)若 $A_{1}, A_{2}, \dots, A_{n}$ 两两互斥, 则
\end{itemize}

\[
P(A_{1} \cup A_{2} \cup \dots \cup A_{n}) = P(A_{1}) + P(A_{2}) + \dots + P(A_{n}).
\]

特别地, 若为对立事件, 则
\[
P(A \cup \overline{A}) = P(A) + P(\overline{A})
\]
即 $P(A) + P(\overline{A}) = 1$, 故 $P(\overline{A}) = 1 - P(A)$.

此节应结合人教 A 版必修二 \textbf{\S{10.1.4} 概率的基本性质} 来学习.

\subsection{频率与概率}

频率 = $\dfrac{\text{频数}}{\text{总数}}$ = $\dfrac{m}{n}$

在相同条件下, 进行大量重复试验, 随机事件 $A$ 发生的频率会在某一个常数附近摆动, 我们将这个常数叫做概率. 通常用频率来估计概率.  

\subsection{事件的独立性}

补充公式:
\[
P(AB) + P(A \overline{B}) = P(A).
\]

由集合论知 $A = A\varOmega = A(B\cup \overline{B}) = (AB)\cup (A\overline{B})$, 而$AB$ 与 $A \overline{B}$ 互斥. 可以结合Venn图理解.

相互独立事件: 若事件 $A$ 是否发生对事件 $B$ 发生的概率没有影响, 则称事件 $A$ 与事件 $B$ 为相互独立事件.

相互独立事件的概率满足
\[
P(AB) = P(A)P(B)
\]

若事件 $A$ 与 $B$ 相互独立, 则 $A$ 与 $\overline{B}$, $\overline{A}$ 与 $B$, $\overline{A}$ 与 $\overline{B}$ 均相互独立.











\end{document}